%%% ---------------------------------------------------------------------------
%%% preamble.tex —— 你自己的宏包与自定义命令
%%%
%%% 主题（common/beamerthemeAcad.sty）已经载入：
%%%   fontspec / unicode-math / acadmath / amsmath / mathtools
%%%   graphicx / booktabs / array / translator / appendixnumberbeamer
%%% 并提供了 evidence 环境、\takeaway、\backuplink、\returnbutton。
%%% **不要载入 enumitem**，它会破坏 \item<2-> 覆盖语法。
%%% ---------------------------------------------------------------------------

\usepackage[
  backend     = biber,
  style       = gb7714-2015,
  maxbibnames = 3,
  minbibnames = 3,
  giveninits  = true,
  doi = false, isbn = false, url = false,
]{biblatex}

%%% ---- 节页 ----
%%% 答辩里节页比会议报告更值得留：它给委员会一个喘息点，也给你一个
%%% 「上一块讲完了」的自然停顿。不想要就注释掉。
\AtBeginSection[]{%
  \begin{frame}[plain, noframenumbering]
    \sectionpage
  \end{frame}%
}

%%% ---- 本答辩专用记号（与论文保持一致）----
\newcommand{\ourmethod}{MSFF-Net}
\newcommand{\ourmethodtwo}{SPR-Head}
\newcommand{\dataset}{DOTA-v2.0}
\newcommand{\featmap}{\mat{F}}
\newcommand{\loss}{\mathcal{L}}
\newcommand{\params}{\vect{\theta}}

%%% ---- 占位图 ----
\newcommand{\placeholder}[3]{%
  \begingroup
    \fboxsep=0pt \fboxrule=0.4pt
    \fbox{\parbox[c][#2][c]{\dimexpr#1-2\fboxrule\relax}{\centering\footnotesize\ttfamily\color{AcadMuted}#3}}%
  \endgroup
}

%%% ---- 演讲备注 ----
%%% 答辩强烈建议写备注：紧张时讲稿是唯一的锚。用 make notes 生成双屏版。
\ifdefined\AcadShowNotes
  \usepackage{pgfpages}
  \setbeameroption{show notes on second screen = right}
\else
  \setbeameroption{hide notes}
\fi
